\begin{afterword}
\summaryhead

The post proposes a holiday, Amazing Breakthrough Day, on April 1st. On that day journalists would report old, understandable scientific breakthroughs as if they had just happened, using April Fools' Day as cover. Sample headlines cover Archimedes explaining why boats float, the impossibility of crossing each of Königsberg's bridges once, and the link between breathing and burning. All are true; they just did not happen yesterday.

The reason, taken from the previous day's post, is that science news covers only the frontier, where findings are often controversial, rest on one experiment, need years of background, and turn out wrong. Readers never see the solid, understandable science on which the frontier is built. Old breakthroughs, like Archimedes's, can be understood without a degree and tested in a bathtub. If an editor objects that readers will not care, the journalist can point to Reddit and Digg, where good explanations of old science get voted up. The post suggests that a new explanation of old science should itself count as news, but calls that too visionary and settles for the holiday as a first step.

\responsehead

Judged as a proposal, the post has a sound diagnosis, a weak plan, and a better idea that it sets aside.

The diagnosis holds up. The post says, with ``often,'' that science news favors new findings that are later shown wrong. A 2017 study of newspaper coverage of biomedical findings found just that: fewer than half of the covered studies were confirmed by later meta-analyses, and journalists favored first findings over later ones.

The plan is to report true old stories on the one day when readers expect headlines to be false. The post says the date gives journalists ``protective cover.'' It does not say what readers are meant to make of a story like ``BOATS EXPLAINED'' on April 1st: whether they are to take it as a joke and later learn it is true, or to learn it from the article. A holiday whose stories may be read as hoaxes is an odd way to teach readers that these stories are true. I found no report that any journalist took the idea up.

The better idea comes near the end, in three sentences: a good new explanation of old science is itself news, and newspapers could treat it that way. That idea needs no holiday. Its only support is the observation that readers of Reddit and Digg vote up such explanations, given without an example, and the post then calls it ``too visionary for a first step.''

The post's main example of an understandable breakthrough mixes two stories. It asks the reader to think of Archimedes's ``Eureka!'' as the moment Archimedes understood why ships float. In the story as told, the bath and the shout concern the volume of a golden crown, and the story itself was first written down two centuries later. Archimedes did explain floating, so the example still serves, but ``events that did, in fact, happen'' is too strong for the bath.

In short: a fair complaint about science news, answered with a holiday whose cover is also a reason to disbelieve its stories, while the better proposal, that explanation is news, gets three sentences.
\end{afterword}
